\section{How the IRT Outcomes Were Constructed}
\label{app:irt}

This appendix explains the measurement procedure in ordinary language. Its
purpose is to place tests from different grades, survey rounds, and years on a
common scale without changing the Year 1 results used in the earlier paper.
Arabic, French, and mathematics remain separate outcomes; there is no claim
that the three subjects measure one common skill.

\subsection{The fixed Year 1 reference}

The published Year 1 analysis supplies both the reference scale and the item
parameters. The latter describe each question's difficulty and how sharply it
distinguishes among students with different levels of achievement. We preserve
the delivered Year 1 student scores and fix the parameters of 104 Arabic items,
85 French items, and 66 mathematics items. The new estimation cannot change
any of those 255 parameter pairs. Automated checks compare the saved values
with the archived files to numerical precision on every run.

We use the Year 1 endline comparison group as the reporting reference. Year 1
baseline is not used because its scores are known to be less reliable. If
\(\mu_s\) and \(\sigma_s\) are the Year 1 endline comparison-group mean and
standard deviation in subject \(s\), the reported score is
\[
 Z_{psn}=\frac{\widehat\theta_{psn}-\mu_s}{\sigma_s}.
\]
The same \(\mu_s\) and \(\sigma_s\) are used for every later grade, round, and
year in that subject. We therefore do not erase changes over time by
standardizing each new test separately.

\subsection{Responses and records included}

Only questions recorded as correct or incorrect enter the IRT models. A value
of 1 is correct, and a value of 0 is incorrect. The Stata code used for the
assessments records ``do not know'' as the tagged missing value \texttt{.a}; we
treat that response as incorrect. Other missing values remain missing. We
exclude continuous measures, counts, oral-reading rates, rubric scores, and
every other non-binary item. None is dichotomized for this analysis.

If a source has duplicate records for the same student and test, the code keeps
the record with the most answered eligible questions. It breaks any remaining
tie using test duration and then source-row order. The number of source rows,
retained rows, and removed duplicates is saved for every assessment group.

\subsection{How one test is linked to another}

An assessment group is a particular year, round, subject, grade, form, and
sample. Different groups often answered different tests. They can be compared
when they share questions whose content and scoring did not change. Such a
question is an \emph{anchor}.

The code first builds a network: assessment groups are the points and accepted
anchors are the links. A group receives a score on the Year 1 scale only when
there is an accepted path from that group to the fixed Year 1 bank. Item maps
and decision logs record the evidence for each occurrence. The project team
has confirmed that, in the operational Year 2 and Year 3 files, an unchanged
item identifier denotes the same question, answer key, scoring, stimulus,
format, and administration, and that an edited question receives a new
identifier. The pipeline records this project decision rather than asking a
statistical procedure to determine whether two identifiers denote the same
question.

\subsection{Allowing questions to behave differently}

Even an unchanged question may work differently in another grade or round.
This is called differential item functioning (DIF). The screening is completed
before treatment effects are estimated. Where treatment status is available,
the preferred screen uses comparison schools; an all-student screen is retained
as a diagnostic. Each candidate anchor is checked with a logistic comparison
and a Mantel--Haenszel comparison among students with similar performance on
the other questions. We adjust for the number of tests with Holm's method and
repeat the screen after removing flagged anchors, for at most three rounds.

The preferred rule examines adjacent-grade differences and change between
baseline and endline. A question is freed only when both the logistic and
Mantel--Haenszel checks indicate DIF and the estimated difference is material.
Freeing means that the implicated administration receives its own item
parameters; the question is not discarded from every test.

Statistical warnings cannot be used mechanically when a test shares only a
small number of questions with the rest of the network. The preferred rule
therefore aims to retain at least five verified anchors on every required link.
After exact item-version mapping, a four-anchor link is allowed only when it is
needed to connect an otherwise isolated assessment group. Questions retained
because of this floor are explicitly flagged in the decision log. The same
rule is applied to Arabic, French, and mathematics.

This safeguard matters for Year 3 Arabic Grades 1 and 2. The baseline Grade 1
and Grade 2 tests share 44 eligible binary questions; Grade 1 shares 41 between
baseline and endline; and the Grade 2 baseline and endline tests share 11. An
earlier, more conservative screen retained only 8, 4, and 2 anchors on the
critical links and consequently left three Arabic assessment groups
disconnected. The missing scores were therefore caused by the screening rule,
not by missing tests, missing binary items, or changed item identifiers. The
preferred consensus-and-connectivity rule restores these groups while keeping
floor-retained questions visible for sensitivity analysis. Table~\ref{tab:mathlinks}
shows the analogous link diagnostics for mathematics, where the network is
also thin in places.

\subsection{The preferred IRT model}

We estimate a two-parameter logistic model. For student \(p\), question \(i\),
and assessment group \(n\),
\[
 \Pr(X_{pin}=1\mid\theta_{pn})
 =\operatorname{logit}^{-1}\!\left[a_i(\theta_{pn}-b_i)\right],
\]
where \(a_i\) is discrimination, \(b_i\) is difficulty, and
\(\theta_{pn}\) is achievement. A student contributes only the questions on
the test actually administered; questions from other forms are missing, not
incorrect.

There is one model per subject. Each model estimates the eligible Year 2 and
Year 3 item parameters together, using all students, while keeping the Year 1
parameters fixed. Estimating the tests together makes use of all links at once
and avoids accumulating a new estimation error at every step in a chain. Once
the item parameters are fixed, each assessment group receives its own estimated
mean and variance. The model therefore does not force different grades or
rounds to have the same achievement distribution.

The preferred person score is the expected value of achievement given the
student's response pattern (EAP). We also calculate a weighted-likelihood
score (WLE) using exactly the same item model. This isolates the choice of
person-score estimator from all anchor and item-parameter decisions. The WLE
routine returned a finite score for every record but issued an iteration-limit
warning for 95 of 56,183 response patterns across 20 assessment groups. We
therefore report WLE as a stress test, not as a competing preferred score. The
preferred EAP scores are unaffected by these warnings.

The freely estimated discrimination parameters have a high numerical upper
bound of 10. This prevents a few nearly separated response patterns from
sending a parameter toward infinity; it is not an item-deletion rule. The cap
does not apply to the fixed Year 1 parameters, which retain their archived
values exactly. The diagnostic table reports how many new parameters reach the
cap, and the alternative measurement methods show whether this stabilization
matters for the treatment estimates. We do not report an unbounded calibration
as though it were a completed robustness result.

The Overall column is not a fourth IRT model. It stacks the three subject
samples after each subject has been placed in Year 1 standard-deviation units.
The subject columns estimate the same regression one subject at a time.

\subsection{Why alternative scale locations do not drive the treatment effect}

A weak link can change the origin of an assessment group's scale even when it
barely changes student rankings. Adding the same constant to every treated and
comparison student in a subject--grade--round group cannot change that group's
treatment-versus-comparison contrast. In the subject-specific regressions,
these common shifts are absorbed by the matched-pair-by-grade time controls.
The Overall estimate also has nearly balanced subject shares across treatment
arms, which is checked in the audit.

Two qualifications matter. First, multiplying a score by a different scale
factor does \emph{not} automatically cancel; it can rescale a treatment effect.
Second, a missing link changes which students can be scored rather than merely
moving the scale origin. We therefore do not rely on cancellation alone. We
re-estimate the treatment effects under alternative links and models, and we
compare IRT with the standardized sum score both on its full sample and on
exactly the students used by the preferred IRT analysis. The robustness tables
show directly whether scale factors or sample coverage matter in practice.

\subsection{Validation and model diagnostics}

The workflow stops if any preferred subject model fails to converge, if a
linked group is not scored, if keys are duplicated, if a non-binary response
enters estimation, if a Year 1 parameter changes, or if the reconstructed
Ministry sum-score regressions fail their tie-out. It also records item-fit
tests that accommodate missing responses, and calculates global fit and
residual item-pair correlations separately within each administered test. These
diagnostics are warnings rather than automatic deletion rules;
Table~\ref{tab:irtdiagnostics} summarizes them. French has the largest share
of item-fit warnings: 120 of 186 tested items. This is a substantive model-fit
caution even though the treatment estimates are stable across measurement
methods; agreement across methods does not by itself establish construct
validity.

The current form registry does not contain complete passage or stimulus
identifiers for every question. Residual item-pair correlations can reveal
possible local dependence, but they are not a substitute for a complete
substantive testlet map. That documentation remains a stated review item rather
than an assumption hidden in the code.

The pipeline also reports 96 unmatched delivered Year 1 baseline score records
in the Cohort 1 one-year panel, 67 in its two-year panel, and 24 in its
three-year panel. These are disclosed as join-coverage warnings. Year 1
endline is the scale reference, and every treatment regression requires an
observed score in both rounds, so unmatched records are excluded rather than
assigned a value.

\subsection{Robustness methods}

Table~\ref{tab:measurement_robustness} keeps the regression fixed and changes
only the construction of the outcome. Its rows cover the choices over which a
researcher could reasonably disagree:

\begin{enumerate}
  \item the preferred EAP score and the WLE score from the same item model;
  \item estimating new item parameters with all students or only comparison
  schools, then scoring both treatment arms;
  \item estimating one model per subject or separate models by grade;
  \item using the preferred anchor rule or a more conservative rule;
  \item estimating all tests together, linking them sequentially, or requiring
  a direct link to Year 1;
  \item applying the targeted DIF procedure without the minimum-link safeguard,
  or allowing a broader set of item parameters to differ;
  \item estimating item parameters in a predetermined subset of schools and
  scoring the remaining schools with those parameters;
  \item using a broader item-match rule in the sequential and direct-link
  approaches; and
  \item using the Ministry's earlier standardized sum score, first on exactly
  the preferred IRT sample and then on every available sum-score observation.
\end{enumerate}

The full-sample Ministry row deliberately preserves the earlier pooled
weighting and fixed-effect implementation so that it can reproduce the tables
already delivered. The preferred IRT and same-sample sum-score rows calculate
equal-cohort weights within the outcome actually being estimated. This prevents
a missing score in one subject from inadvertently changing that cohort's
influence in another subject.

The companion sample table reports the number of students retained by each
method relative to the preferred IRT analysis. A dash means that a method does
not connect every group needed for that comparison or that the prespecified
regression cannot be estimated. The code never substitutes a different cohort,
drops a required regression term, or calls an estimate based on fewer than
three subjects an Overall effect.

The preferred decision file contains 266 item-by-link rows whose parameters
are allowed to differ, 39 rows retained to preserve a required link, and 2,868
rows with a material difference that did not also satisfy the two-test
statistical rule. These are administration-specific decisions, not counts of
globally defective questions.

We also registered two deliberately demanding models. A model that estimates a
separate group distribution during joint calibration remains an unfinished,
quarantined extension; no result from it is reported. A joint model based only
on the broader item-ID rule
converged under the declared numerical bound, but remains quarantined because
that broader matching rule is not accepted anchor evidence. Both attempts
remain in the method registry but are not treated as competing preferred estimates.
A Rasch model was not treated as an
interchangeable robustness outcome because it would replace the fixed Year 1
two-parameter scale with a different measurement model and would require its
own defensible link. Treatment-related DIF is also reserved for a later
sensitivity analysis; treatment effects never determine the preferred anchors.

Finally, the automated field \texttt{final\_outcome\_approved} remains zero.
External release still requires confirmation of canonical source choices,
completion of the Arabic and French occurrence review, and complete
stimulus/testlet metadata. These gates concern scientific sign-off; they do not
indicate a failed estimation or a numerical error in the treatment effects.

\clearpage
\begin{landscape}
\input{tables/08_measurement_robustness.tex}
\end{landscape}

\clearpage
\begin{landscape}
\input{tables/09_measurement_sample_retention.tex}
\end{landscape}

\clearpage
\input{tables/10_math_link_diagnostics.tex}

\clearpage
\input{tables/11_primary_irt_diagnostics.tex}
